\chapter{A stochastic laboratory for temporal entropy inference}
\label{ch:vl-thermodynamic}

This laboratory generates hidden trajectories, noisy detector outputs and
experimental timing before analyzing the observed counts. Its purpose is to
test the full inference chain under declared apparatus models. It supplies
simulation evidence; none of the following records came from a physical
laboratory. The time unit and transition rates are illustrative.

The experiment specification, independent pilot and held-out seeds,
record budgets, detector laws, confidence radii and original inference
scores were saved in \texttt{contracts/thermodynamic\_lab.json} before
production. Twenty pilot datasets and 200 held-out datasets per case use
disjoint random streams. Pilot outcomes remain in the release. The
subsequent comparison with a direct, predeclared curvature contrast uses
the same held-out records and is reported separately. This separation of
design, data generation and evaluation follows established simulation-study
practice \cite{vlMorris2019}.

\section{Trajectory and apparatus construction}

The main physical model is a ring with $N=40$ states and translation-invariant
forward and backward rates $a,b$. Each independent pair record begins with
$X_0$ uniform on the ring, which is its exact stationary distribution. For
the selected duration $U$, generate independent counts
\begin{equation}
 N_+\sim\operatorname{Poisson}(aU),\qquad
 N_-\sim\operatorname{Poisson}(bU),\qquad
 X_U=(X_0+N_+-N_-)\bmod40.
\end{equation}
This samples the continuous-time random walk exactly: independent forward
and backward Poisson clocks generate its jumps. The simulator does not draw
matching bits from an analytically supplied endpoint probability.

A fixed detector has conditional mean $f_j\in[-1,1]$. Independent readouts
are generated at both physical endpoints from
\[
 \Pr(Y=+1\mid X=j)=\frac{1+g f_j}{2},\qquad 0\le g\le1.
\]
Thus a gain $g$ multiplies the stationary two-time correlation by $g^2$.
The first menu uses a deterministic half-ring detector, $f_j=+1$ on 20
consecutive states and $-1$ on the others. The second uses the fixed rational
approximation to a cosine detector specified in the frozen contract.
These are declared detector models, not empirical calibrations.

The revival menu observes lags $1.6$ and $2.1$, with 5,000 independent
records per lag. The curvature menu observes $0.15$, $0.30$ and $0.45$,
with 50,000 records per lag. They are separate experiments with different
detectors and observation windows; their results are not presented as a
matched efficiency comparison. Within each menu, every inference procedure
receives exactly the same records.

The driven ring has $(a,b)=(20,1)$ and stationary entropy-production rate
$19\log20\simeq56.9189$ in natural entropy units per model time. The
equilibrium control has $(a,b)=(10.5,10.5)$, preserving total jump activity
while setting stationary entropy production to zero. A noisy detector
control uses $g=0.2$, equivalent to an independent bit flip with probability
$0.4$ after each ideal readout.

An independent timing control replaces each nominal lag $t_k$ by
$U_k=t_kV$, where $V$ is uniform on $[0.7,1.3]$ and independent of physical
paths and detector noise. The same scale law is used at every lag. For a
reversible process at $t_k=k\Delta$,
\[
 C(U_k)=m^2+\sum_\ell c_\ell\mathbb E
       [e^{-\lambda_\ell k\Delta V}],\qquad c_\ell\ge0.
\]
After averaging over $V$, this is still a nonnegative moment mixture with
$x=e^{-\lambda_\ell\Delta V}\in[0,1]$. Consequently the reversible
correlation cone is preserved, and the expected duration remains nominal
because $\mathbb E[V]=1$. Arbitrary setting-dependent timing errors do not
inherit this conclusion.

\section{Blind inference and matched accounting}

Inference receives the number of matching endpoint pairs at each lag,
not hidden states or jump counts. With $K$ settings, $n$ records at each,
and empirical matching probabilities $\widehat p_k$, the simultaneous box
\begin{equation}
 p_k\in[\max(0,\widehat p_k-r),\min(1,\widehat p_k+r)]
 \quad\text{for all }k
\end{equation}
has failure probability at most $2K e^{-2nr^2}$. The chosen radii are
$r=0.021$ for revival and $r=0.007$ for curvature. Rational logarithm
enclosures verify that both failure bounds are at most $0.05$.

Three original procedures use this identical simultaneous box. A
negative-correlation test rejects when an upper probability endpoint is
below $1/2$. A revival test rejects when the upper endpoint at an earlier
lag is below the lower endpoint at a later lag. A fixed dual witness
evaluates the previously specified positive likelihood scores against the
full reversible moment cone. Its exact polynomial normalizer is checked
once, and each dataset is evaluated using the same rational logarithm
enclosures as the original certificate. The latter is a fixed witness from
the full cone; it is not an optimization over every valid witness.

These tests need no correction merely because they are compared: each
rejection is impossible on the same joint confidence event under its
reversible null. Accordingly even the union of their null rejections is
bounded by the box failure probability. This does not extend to changes in
the measurement contract.

The revival menu uses 10,000 pair records and nominal observation exposure
$5000(1.6+2.1)=18,500$ per dataset. The curvature menu uses 150,000 pair
records and exposure $50,000(0.15+0.30+0.45)=45,000$. Actual elapsed exposure
is stored for every dataset, including jitter. Physical preparation,
stationary resetting, detector operation and data processing costs are not
calibrated or included in these exposure totals. The calculation does not
equate wall-clock runtime with thermodynamic energy.

\section{Held-out results and a useful failure}

Table~\ref{tab:vl-thermodynamic} gives rejection counts in 200 independently
generated held-out datasets per case. All entries are simulated outcomes.
The JSON report also gives exact-binomial Clopper--Pearson confidence
intervals for each Monte Carlo rejection probability
\cite{vlClopper1934}, evaluated numerically. For example, $0/200$ gives an
upper 95\% endpoint of approximately $0.0183$, while $200/200$ gives a lower
endpoint of approximately $0.9817$. Neither outcome establishes zero or
unit rejection probability.

\begin{table}[ht]
\centering
\small
\begin{tabular}{llrrr}
\hline
Menu & Simulated case & Negative & Revival & Fixed dual\\
\hline
Curvature & Driven & 0 & 0 & 1\\
 & Equilibrium & 0 & 0 & 0\\
 & High detector noise & 0 & 0 & 0\\
 & Independent scale jitter & 0 & 0 & 0\\
 & Detector memory$^*$ & 0 & 0 & 1\\
Revival & Driven & 0 & 200 & 200\\
 & Equilibrium & 0 & 0 & 0\\
 & High detector noise & 1 & 0 & 0\\
 & Independent scale jitter & 2 & 200 & 0\\
 & Detector memory$^*$ & 0 & 200 & 200\\
 & Prepared equilibrium$^*$ & 200 & 0 & 0\\
\hline
\end{tabular}
\caption{Rejections among 200 held-out simulated datasets, with matched
records and exposure within each menu. Asterisks identify controls outside
the stationary, independently read-out detector contract. Their rejections
must not be interpreted as dissipation of the simulated physical sample.
The fixed dual uses the original predeclared scores, not all possible
moment-cone witnesses.}
\label{tab:vl-thermodynamic}
\end{table}

The favorable revival experiment works with the planned record budget,
and equilibrium controls produce no rejections in these runs. Independent
readout noise substantially suppresses useful information although the
sample's hidden stationary entropy production is unchanged.

The curvature experiment exposes a practical limitation. Its model curve
has the required departure from convex equilibrium decay, but the
combination of coordinate confidence intervals and the original score
choice removes most of this witness's usable separation at this budget.
The fixed witness rejects only once in 200 held-out datasets.
This is not a successful practical detection design at that record count.
It does not establish that optimized moment-cone inference has low power.
The subsequent matched comparison includes both safely optimized box
witnesses and direct-contrast inference, which concentrates uncertainty
in a predeclared informative direction.

Timing jitter reveals a second distinction between validity and power.
Common random scaling preserves the reversible cone, yet the original
likelihood scores are poorly positioned for the shifted matching
probabilities: the fixed dual never rejects, while the broader revival
inequality rejects in all 200 datasets. A cone can remain the correct
null class while one selected witness becomes ineffective. Simultaneous
confidence boxes permit safe subsequent witness optimization, but these
frozen-score results are retained without replacing the failed regime.

\section{Controls that defeat an unqualified interpretation}

The memory control evolves the sample with the equilibrium ring generator,
but the detector stores its first bit and repeats it with a
setting-dependent probability specified in the contract. This deliberately
violates the assumption that each readout is a fixed channel of the current
physical state with independent noise. The resulting revival record
imitates the driven record and is rejected in all 200 datasets by both
revival and the fixed dual. Sampling confidence about those observed
probabilities cannot establish the missing detector contract. The
sample's generator has zero stationary entropy production; the separate
memory device's physical costs have not been modeled.

The preparation control is a detailed-balanced two-state process with
rates $0\to1$ equal to $9$ and $1\to0$ equal to $1$. Direct Gillespie
simulation starts every record in state $0$, rather than in its stationary
law $(0.1,0.9)$. The matching probability is
$0.1+0.9e^{-10t}$, and the negative-correlation test rejects every held-out
dataset. The stationary dissipation interpretation is inapplicable.
The preparation-relative entropy is $D(\delta_0\Vert\pi)=\log10$;
the earlier preparation-or-housekeeping theorem explicitly allows this
source of apparent temporal separation.

\section{Reproducibility and remaining physical work}

The release retains the full 150,000-record first held-out curvature dataset
in a compressed archive: initial and final hidden states, forward and
backward jump counts, elapsed times, both detector bits and setting labels.
These hidden arrays support auditing only. Separate compact archives retain
counts and exposure for every pilot and held-out dataset. The frozen
contract and deterministic archives have recorded SHA-256 digests.

An independent matrix-exponential calculation checks the mean matching
probabilities. For jitter, 32- and 64-node quadratures must agree within
$10^{-11}$. Reports include empirical joint confidence-box coverage against
these numerical oracles: all 11 cases cover the model probability vector
in between 198 and 200 of their 200 held-out datasets. This includes the
outside-contract controls, emphasizing that correct statistical coverage
does not establish the physical interpretation. Tests also compare the direct Gillespie control
with its independent analytic transition law and compare compiled inference
with the original exact certificate implementation.

This adds an executable apparatus stress test to existing thermodynamic
inference research. Finite resolution and finite statistics have already
been studied explicitly \cite{vlFritz2025}; entropy inference from
time-symmetric observations is also established \cite{vlSkinner2021}.
The comparison here is deliberately confined to the shared endpoint record.
Methods requiring transition labels, waiting times or resolved currents
have different access and are not ranked as inferior from this experiment.

Finite-duration detector integration, imperfect physical resets,
unmodeled sample drift and independently calibrated hardware energy remain
experimental dependencies. The point-readout simulator and independent
stationary preparation do not validate them.
